
\subsection{Adversarial Search}

The adversarial search agent is designed for decision-time planning under
imperfect information. Unlike value-based agents, it does not learn parameters
from self-play. Instead, it asks the following local question at every move:
given the current information set, which legal action is most robust against
likely responses from the other two players?

The main difficulty is that Fight the Landlord is not a perfect-information
game. The current player observes only an information set

\[
    I = (h_i, H, m_{\mathrm{last}}, c, r_i),
\]

where \(h_i\) is the player's own hand, \(H\) is the public action history,
\(m_{\mathrm{last}}\) is the last valid move, \(c\) records the remaining card
counts of the three players, and \(r_i\) is the role of the current player. The
hidden hands of the other two players are unknown. Therefore, a direct minimax
search over the true state is impossible.

To approximate the hidden state, we use determinization. Let
\(\Omega(I)\) denote the set of all complete states consistent with the
information set. The agent samples \(N\) possible complete states
\(s_1,\ldots,s_N \sim \Omega(I)\) by randomly assigning the unknown cards to
the two unobserved hands according to their remaining card counts. For a root
action \(a\), the sampled search value is

\[
    Q_{\mathrm{search}}(I,a)
    =
    b(I,a)
    +
    \frac{1}{N}
    \sum_{k=1}^{N}
    V_d(T(s_k,a)),
\]

where \(b(I,a)\) is an immediate tactical bonus or penalty, \(T\) is the state
transition after playing \(a\), and \(V_d\) is a depth-limited search value. The
term \(b(I,a)\) handles short-term tactical cases such as finishing the hand,
blocking a dangerous enemy, releasing a teammate, avoiding unnecessary bombs,
and penalizing risky high singles.

Inside each determinized state, the search is role-aware. If the current node
belongs to the root agent, the agent assumes it can choose the best continuation:

\[
    V(s) = \max_{a \in A_s} V(T(s,a)).
\]

For teammate nodes, using a pure maximum is too optimistic because the teammate
does not know the root player's hidden cards and may not follow the searched
line perfectly. We therefore combine a policy expectation with a small maximum
component:

\begin{align*}
    V_{\mathrm{team}}(s)
    = & \ (1-\mu)
    \sum_{a \in A_s}
    \pi_{\tau}(a \mid s) V(T(s,a))        \\
      & + \mu \max_{a \in A_s} V(T(s,a)).
\end{align*}

Here \(\pi_{\tau}\) is a softmax policy over hand-crafted action scores, and
\(\mu\) controls how much we trust the teammate to find the best cooperative
line.

For enemy nodes, a pure minimax update would be too pessimistic. It would assume
that the opponent knows the determinized full state and always selects the most
damaging action. However, using only a soft policy can also be too optimistic,
because it may underweight the opponent's strongest punishments. We therefore
use a symmetric mixed model: most of the value comes from a bounded-rational
policy expectation, while a small component still accounts for the worst
adversarial response:

\begin{align*}
    V_{\mathrm{enemy}}(s)
    = & \ (1-\nu)
    \sum_{a \in A_s}
    \pi_{\tau}^{\mathrm{enemy}}(a \mid s)
    V(T(s,a))                             \\
      & + \nu \min_{a \in A_s} V(T(s,a)).
\end{align*}

The action probabilities are computed from heuristic action scores:

\[
    \pi_{\tau}(a \mid s)
    =
    \frac{\exp(g(s,a)/\tau)}
    {\sum_{a'} \exp(g(s,a')/\tau)}.
\]

Here \(\nu\) controls how much pessimism is retained in the opponent model.
This model is less paranoid than minimax but still biased toward strong
responses. In practice, the search also prunes actions to a small top-\(K\)
candidate set at both the root and internal nodes, which keeps the computation
tractable.

When the search reaches a terminal state, the value is a large positive or
negative constant depending on whether the root team wins:

\[
    V(s_{\mathrm{terminal}})
    =
    \begin{cases}
        +R_{\mathrm{win}}, & \text{if the root team wins}, \\
        -R_{\mathrm{win}}, & \text{otherwise}.
    \end{cases}
\]

At non-terminal leaf nodes, the agent falls back to a hand-crafted state
evaluation. This evaluation rewards the root team for having fewer and better
structured remaining cards, penalizes dangerous enemies with very few cards
left, and adds bonuses for initiative, last valid control, bombs, and high
control cards.

Finally, the agent does not blindly trust the sampled search result. Because
determinization and shallow search can be noisy, the chosen action must be
sufficiently better than the baseline greedy candidate or the second-best
candidate:

\[
    a_{\mathrm{search}} = \arg\max_a Q_{\mathrm{search}}(I,a).
\]

The search action is accepted only if

\begin{align*}
    Q(a_{\mathrm{search}}) - Q(a_{\mathrm{base}})
     & \ge \delta_{\mathrm{base}},   \\
    \text{or}\quad
    Q(a_{\mathrm{search}}) - Q(a_{\mathrm{second}})
     & \ge \delta_{\mathrm{second}}.
\end{align*}

Otherwise, the agent falls back to the baseline action. This gate makes the
agent more stable: adversarial search is used when it gives strong evidence for
an improvement, but the policy avoids being dominated by sampling noise or
incorrect hidden-card assumptions.

Overall, this method can be viewed as a practical approximation to
imperfect-information minimax: hidden cards are sampled, a shallow team-aware
search is performed on each sampled complete state, and a protected decision
rule combines the search result with a strong heuristic baseline.
